\section{Results}

\subsection{Publication cohort and response quality}

The cleared publication pool comprised 40 participants (33 women, 7 men; mean age 23.2 years, SD 4.0, range 19--37). Twenty-one completed Control first and 19 completed Viscereality first. Across four acquisition versions, 16{,}298 of 16{,}320 planned trial slots were observed. One admitted partial session lacked 22 planned trials but retained all 20 Condition $\times$ Phase $\times$ SOA cells; analyses excluding this participant are reported below. The pool included 9{,}588 observed Audio--Tactile trials, 4{,}793 Baseline trials, and 1{,}917 Catch trials. Audio--Tactile and Baseline hit rates were 91.8\% and 90.8\%, respectively; the Catch false-alarm rate was 1.2\%. No participant exceeded the preregistered 30\% Catch false-alarm exclusion criterion.

Of 13{,}153 Audio--Tactile and Baseline trials labelled as hits, 105 (0.8\%) had no recorded reaction time: 80 Audio--Tactile and 25 Baseline. All 105 came from the later three acquisition versions; they remained hits in response-quality summaries but could not enter reaction-time models. Of the remaining 13{,}048 hits with a recorded reaction time, 363 were removed by the preregistered participant-wide mean $\pm$ 2.5-SD rule. The resulting reaction-time dataset contained 8{,}495 Audio--Tactile and 4{,}190 Baseline trials. The 200-ms lower cutoff used in an earlier analysis of the first 16 participants was not specified in the preregistration and was not applied here; adding it removes one further hit without changing the sigmoid-viability conclusion.

\subsection{Preregistered \texorpdfstring{Condition $\times$ Phase $\times$ SOA}{Condition x Phase x SOA} test}

The primary analysis compared nested maximum-likelihood linear mixed models of individual Audio--Tactile reaction times, with participant random intercepts and all lower-order Condition, Phase, and numeric-SOA terms in both models. Adding the preregistered Condition $\times$ Phase $\times$ SOA interaction did not improve fit, $\chi^{2}(1)=0.006$, $p=.939$. With Condition and Phase coded symmetrically and SOA standardized, the interaction estimate was 0.42 ms per SOA standard deviation (SE 5.47; approximate 95\% Wald CI $[-10.30, 11.13]$). The sampled 300--2700-ms span was 2.73 SOA standard deviations, corresponding to a far-to-near three-way contrast of about 1.1~ms with an approximate 95\% interval of $[-28.2, 30.4]$~ms. Thus the data did not detect the predicted Viscereality-specific change in the respiratory-phase SOA gradient. This interval does not establish that the effect is exactly zero.

The conclusion was stable when SOA was treated as a five-level factor, $\chi^{2}(4)=1.94$, $p=.747$, and when the documented partial session was excluded (numeric-SOA $p=.740$; categorical-SOA $p=.827$; $n=39$). The registered secondary Control-only Phase $\times$ SOA comparison was also not detected (numeric-SOA $p=.420$; categorical-SOA $p=.513$). A preregistration-cleaned 2 $\times$ 2 $\times$ 5 repeated-measures analysis of participant cell means likewise found no three-way interaction, $F(4,156)=0.82$, Greenhouse--Geisser-corrected $p=.494$. A trial regression with participant-cluster-robust standard errors estimated the numeric three-way contrast at $-2.18$ ms per SOA SD, 95\% CI $[-10.79, 6.43]$, $p=.611$. An equal-participant analysis of five-SOA slopes gave $-0.42$ ms per SOA SD, 95\% CI $[-8.67, 7.82]$, $p=.918$ (sign-flip $p=.921$). A mixed model with a participant-specific SOA slope also found no three-way interaction, $\chi^{2}(1)=0.010$, $p=.921$. Retained Audio--Tactile RTs were right-skewed (sample skewness 1.28); a log-RT trial regression with participant-cluster-robust uncertainty likewise did not detect the interaction ($p=.520$). Subtracting Baseline RT at each same nominal SOA label gave an exploratory participant-level interaction of 1.45 ms per SOA SD, 95\% CI $[-20.40, 23.29]$, $p=.894$; this is not a verified cue-time-matched comparison. These are sensitivity analyses rather than replacements for the registered mixed model. The exploratory one-sample JZS test of the raw Audio--Tactile participant slope contrasts with Cauchy scale $\sqrt{2}/2$ yielded $\mathrm{BF}_{10}=0.171$ ($\mathrm{BF}_{01}=5.83$); it is not the preregistered mixed-model Bayes factor, whose full prior and model specification were not stated. The participant contrast distribution departed from normality (Shapiro--Wilk $p=.007$), limiting the JZS model interpretation.

Audio--Tactile hit rates ranged from 91.0\% to 93.0\% across SOAs. An exploratory participant-cluster-robust binomial model did not detect a Condition $\times$ Phase $\times$ SOA interaction in hit probability (log-odds estimate $-0.033$ per SOA SD, 95\% CI $[-0.277, 0.210]$, $p=.785$). Mean within-condition trial position was similar across SOAs (101.8--103.3 of 204 planned slots). Adjusting the cluster-robust RT regression for linear within-condition trial position and Condition $\times$ Phase-specific time trends left the three-way estimate at $-2.22$ ms per SOA SD, 95\% CI $[-10.80, 6.37]$, $p=.604$. In the post hoc missing-hit-RT stress test, the primary interaction estimate ranged from $-4.86$ to $5.92$ ms per SOA SD under the opposing 200-ms shifts (likelihood-ratio $p=.376$ and $p=.281$, respectively); imputing the within-cell median alone gave $p=.975$. These fixed scenarios do not identify the reaction times of missed trials, recover missing hit RTs, or bound arbitrary departures from the observed RT distribution.

\begin{figure}[htbp]
\centering
\includegraphics[width=\textwidth]{study5_current_rt_soa_curves.png}
\caption{\textbf{Audio--Tactile reaction time by SOA, Condition, and instructed respiratory phase in the 40-person publication pool.} Each point is the mean of participant reaction-time means at that SOA; each participant has equal weight. All 20 plotted cells contain 40 participants and 413--442 retained reaction times. Error bars are 95\% $t$ intervals across participant means and show between-person uncertainty for each plotted mean. Longer SOAs correspond to nearer auditory source positions. The interaction inference comes from the trial-level mixed model, not from overlap of these intervals.}
\label{study5:fig:current_rt_soa}
\end{figure}

\subsection{Sigmoid viability and descriptive PPS shape}

Under the historical trial-level implementation of the preregistered four-parameter sigmoid fit, none of the 160 participant $\times$ Condition $\times$ Phase cells reached $R^{2}\geq .70$; consequently no participant supplied four qualifying curves. Sigmoids had lower in-sample RMSE than linear fits in 156 of 160 cells, as specified for comparison in the registration, but the median improvement was only 0.87~ms against a median linear RMSE of 110.81~ms. Median sigmoid $R^{2}$ was .038 (maximum .419), and 54 fits touched a parameter bound. The unpenalised same-data RMSE comparison does not establish a stable sigmoid. If individual trials were first averaged to five SOA means, 73 of 160 cells reached the threshold and two participants passed in all four cells, but neither had all four fits free of parameter bounds. Four parameters fitted to five SOA means leave only one residual degree of freedom. The four participant-weighted group curves can be displayed descriptively, but all four unsmoothed group sigmoid fits reached at least one parameter bound. We therefore did not test participant midpoint or sharpness contrasts as interpretable PPS-boundary outcomes. The registered exploratory MAIA-2 moderation and pictographic associations with boundary modulation were likewise not tested because they require interpretable participant boundary estimates. No smoothing or post hoc parameter constraints were substituted for this gate.

Replacing the registration's approximate distance labels with positions implied by its stated linear trajectory left the trial-level gate at 0/160 cells and the five-mean gate at 73/160 cells, with the same two participants passing all four mean-fit cells but neither free of bound-hitting curves. Bound-hitting five-mean fits increased from 88 to 102 of 160. Thus the gate conclusion is unchanged, although individual shape estimates depend on the coordinate convention.

Audio--Tactile minus condition-and-phase-averaged tactile-only Baseline RT was $-17.94$ ms, 95\% CI $[-26.44, -9.44]$, $p<.001$; trial-relative cue timing differed between modalities in the verified source subset. The exploratory far-minus-near Audio--Tactile RT contrast was 16.15 ms, 95\% CI $[0.28, 32.03]$ (without the partial session: 15.22 ms, 95\% CI $[-0.96, 31.41]$). Baseline RT changed by 13.00 ms across its nominal delay labels, 95\% CI $[-6.42, 32.43]$. The difference between these two label-based contrasts was 3.15 ms, 95\% CI $[-18.11, 24.40]$ (without the partial session: 1.49 ms, 95\% CI $[-20.06, 23.05]$). Because equal labels did not mean equal cue times in the verified source subset, this difference does not remove temporal-expectancy effects. The raw SOA gradient neither isolates auditory proximity nor identifies a PPS boundary.

All 167 available original block recordings from 14 cleared participants matched their played-order references, yielding 5010 tactile cues. In every block, the median Audio--Tactile cue offset exceeded the Baseline offset at the same nominal label by 500~ms, with no within-trial-type deviation at the 10-ms detection resolution. A separate cleared source session with surviving event times showed the same 500-ms difference in 12/12 condition-by-block median comparisons; the pattern held within 1~ms for 339/341 cue trials, with two individual timing anomalies. For 201 Audio--Tactile trials in that session with both detected onsets, the measured auditory-to-tactile interval differed from its nominal SOA by a median absolute 0.3~ms. Direct timing evidence therefore covers 15/40 cleared participants, all from the earliest acquisition version; the other 25 remain unverified.

\subsection{Baseline assumptions and acquisition sensitivity}

Tactile-only Baseline reaction time was 19.83 ms longer after inhalation than after exhalation, 95\% CI $[4.51, 35.16]$, unadjusted $p=.013$. The exploratory Baseline Condition $\times$ Phase contrast was 17.14 ms, 95\% CI $[0.45, 33.83]$, unadjusted $p=.044$; the latter is uncertain after accounting for multiple assumption checks. The registered expectation of phase-invariant Baseline responses was therefore not supported. Raw Audio--Tactile reaction time was retained as the primary outcome, and Baseline-subtracted curves were treated as descriptive. At the farthest auditory SOA, Audio--Tactile minus Baseline RT was $-8.06$ ms, 95\% CI $[-22.63, 6.50]$, using the historical condition-and-phase reference. Pairing Baseline by the same nominal SOA label instead gave $-22.85$ ms, 95\% CI $[-38.94, -6.76]$. Neither reference is established as a cue-time-matched test of the registered no-far-distance-facilitation assumption. For the two sound types available across the whole pool, the participant-level pink-minus-white Audio--Tactile difference was $-4.02$ ms, 95\% CI $[-10.89, 2.86]$; this does not establish equivalence.

Exploratory paired-contrast JZS Bayes factors with Cauchy width $\sqrt{2}/2$ were $\mathrm{BF}_{10}=0.261$ for the Baseline Condition difference, 3.362 for the Baseline Phase difference, 0.305 for far Audio--Tactile versus the condition-and-phase mean Baseline, 5.864 for far Audio--Tactile versus Baseline at the same nominal SOA label, and 0.326 for pink versus white noise. The Baseline Phase participant contrast departed from normality (Shapiro--Wilk $p=.021$). These model-dependent factors are exploratory implementations of the registration's incompletely specified Bayesian tests; the same-label factor is not a matched-timing comparison, and the pink-white result does not establish equivalence across all four planned sound types.

Within the earliest acquisition version, all 17 cleared participants had Audio--Tactile trials with all four planned sound types. Equal-participant mean RTs were 737.35~ms for blue, 717.23~ms for brown, 715.37~ms for pink, and 721.37~ms for white noise; an exploratory Friedman comparison gave $\chi^{2}(3)=7.38$, $p=.061$. This within-version result does not establish equivalent RTs across sounds or generalise to later versions with fewer sound types.

\begin{figure}[htbp]
\centering
\includegraphics[width=\textwidth]{study5_delay_matched_facilitation.png}
\caption{\textbf{Exploratory Audio--Tactile minus tactile-only Baseline RT paired by nominal SOA label.} Each point is the mean of 40 participant differences after averaging within participant $\times$ Condition $\times$ Phase $\times$ nominal label $\times$ trial type. Negative values indicate faster Audio--Tactile responses. Baseline trials had no sound; source timing checks in 15 participants found their tactile cues 500~ms earlier within the trial than Audio--Tactile cues with the same label. Thus the curves are not verified cue-time-matched facilitation. Error bars are pointwise 95\% $t$ intervals across participants; lines are guides to the eye, not fitted sigmoids.}
\label{study5:fig:delay_matched_facilitation}
\end{figure}

Condition-order subgroups yielded three-way slope contrasts of $-4.19$ ms per SOA SD (Control first; $n=21$, 95\% CI $[-14.64, 6.27]$) and $3.74$ ms (Viscereality first; $n=19$, 95\% CI $[-10.05, 17.53]$). These estimates are imprecise, and order is partly confounded with acquisition version. The four acquisition versions also differ in available noise types and in documented reaction-time onset definitions. Excluding two participants with explicitly different RT onset definitions left the participant-cluster-robust three-way estimate at $-1.38$ ms per SOA SD, 95\% CI $[-10.37, 7.60]$, $p=.757$ ($n=38$). The remaining source-recorded reaction times were still not uniformly remeasured from original waveforms. The pooled analysis thus supports a null detection of the preregistered interaction under the available recorded timing data, with these provenance limits carried into interpretation.
